\chapter*{ВВЕДЕНИЕ}
\addcontentsline{toc}{chapter}{ВВЕДЕНИЕ}
\sloppy

Большие языковые модели (LLM) всё активнее применяются для~автоматизации
рутинных операций при~проектировании образовательных материалов.
Мировой рынок корпоративного обучения, по~оценкам аналитических агентств,
превысит \$51~млрд к~2030~году~\cite{edtech_market_report_2024}.
Главный сдерживающий фактор роста~--- высокая стоимость
и~длительность создания качественного учебного контента.
Средний цикл разработки одного корпоративного курса силами методиста-эксперта
составляет от~нескольких дней до~нескольких недель и~требует педагогических
компетенций, которыми не~всегда обладают эксперты предметной области.

Существующие ИИ-инструменты для~ускоренной подготовки курсов (Coursebox~AI, iSpring~AI)
предоставляют одношаговую генерацию на~основе единственного промпта и,
как~правило, не~поддерживают научно обоснованные методологии проектирования учебного процесса
(ADDIE, SAM, Backward~Design).
Мультиагентные LLM-конвейеры для~решения этой задачи
остаются малоизученными: отсутствуют систематические сравнения архитектур
на~воспроизводимом российском корпусе документов и~внешнем зарубежном сравнительном контуре, а~также воспроизводимые бенчмарки
с~привязкой не~только к~таксономии Блума, но и~к~согласованности курса,
качеству онлайн-курса и~переносу обучения в~практику.

\textbf{Объект исследования}~--- образовательная платформа Adapstory~(AI~LMS),
предоставляющая инструменты создания и~управления учебными курсами
для~B2B-клиентов.

\textbf{Предмет исследования}~--- микросервис ИИ-методист:
интеллектуальный проактивный методический ассистент, поддерживающий
методиста, эксперта и~разработчика курса при~создании, упаковке и~актуализации
образовательных продуктов на~основе LLM с~применением мультиагентной
оркестрации.

\textbf{Цель работы}~--- разработка и~валидация интеллектуального
методического ассистента, сокращающего рутинную нагрузку при~подготовке
структуры и~черновых материалов учебных курсов на~основе входных
документов заказчика при~сохранении за~человеком методического решения
и~контроля качества, измеримого по~таксономии Блума.

Для~достижения цели поставлены следующие задачи:
\begin{enumerate}
  \item провести анализ предметной области: Customer~Development с~B2B-аудиторией,
    продуктовый обзор конкурентов и~систематический обзор научных подходов;
  \item оценить рыночный потенциал решения (TAM/SAM/SOM) и~сформулировать
    исследовательские гипотезы;
  \item спроектировать мультиагентную архитектуру конвейера поддержки
    проектирования курсов
    и~обосновать выбор технологического стека через воспроизводимые эксперименты;
  \item реализовать систему: контур обработки документов (Document~Ingestion~Pipeline, Haystack~2.x),
    граф агентов LangGraph, LLM~Gateway, FastAPI-микросервис;
  \item провести валидацию системы: автоматизированную оценку
    (Ragas, покрытие уровней таксономии Блума),
    экспертный слепой тест (n~=~5--8~методистов) и~пилотное исследование с~B2B-клиентами;
  \item оценить бизнес-потенциал: unit-экономику, финансовую модель
    и~провести Customer~Development для~подтверждения соответствия решения выявленной проблеме.
\end{enumerate}

\textbf{Научная новизна} работы состоит в~следующем:
предложена и~экспериментально верифицирована мультиагентная архитектура
поддержки проектирования учебных курсов, впервые сочетающая
(1)~вариативную методологическую модель проектирования (ADDIE / SAM / Backward~Design)
как~переключаемый системный промпт,
(2)~графовую оркестрацию с~сохраняемым состоянием и~участием человека
(LangGraph~StateGraph),
и~(3)~RAG-слой на~воспроизводимом российском корпусе открытых
и~официальных документов с~отдельным зарубежным сравнительным контуром
(Qdrant~+~deepvk/USER-bge-m3).
Архитектура превосходит одношаговую базовую конфигурацию
по~достоверности относительно источника и~покрытию уровней
таксономии Блума.

\textbf{Положения, выносимые на защиту:}
\begin{enumerate}
  \item Мультиагентный конвейер с~вариативной методологической моделью
    (ADDIE/SAM/Backward~Design) сокращает время подготовки первого
    методического черновика курса с~X~часов до~Y~минут
    при~сохранении покрытия уровней таксономии Блума $\geq 70\%$.
  \item Графовая архитектура с~сохраняемым состоянием (LangGraph)
    обеспечивает статистически значимо более высокое качество структуры курса
    по~сравнению с~одношаговым подходом на~той же базовой LLM.
  \item Система прошла этап подтверждения соответствия решения проблеме:
    $N$~B2B-пилотов
    подтверждены письмами о~готовности от~ЛПР.
\end{enumerate}

Текущее состояние проекта: разработан и~протестирован MVP,
включающий контур обработки документов и~агентный конвейер
поддержки проектирования курса с~методологией~ADDIE; проведены
B2B-интервью и~получены LOI-письма от~корпоративных заказчиков.

Работа состоит из~введения,
\ifomitvalidationchapter
трёх глав, заключения,
\else
четырёх глав, заключения,
\fi
раздела об~использовании ИИ, списка использованных источников
и~десяти приложений.
Глава~1 содержит анализ предметной области и~постановку задачи.
Глава~2 описывает дизайн эксперимента и~обоснование выбора технологий.
Глава~3 посвящена технической реализации системы.
\ifomitvalidationchapter
\else
Глава~4 содержит результаты валидации, анализ ограничений
и~бизнес-оценку.
\fi
